\documentclass[11pt]{article}
\usepackage[margin=1in]{geometry}
\usepackage{amsmath,amssymb,amsthm}
\usepackage{booktabs}
\usepackage{array}
\usepackage[hidelinks]{hyperref}

\newtheorem{theorem}{Theorem}
\newtheorem{proposition}{Proposition}
\newtheorem{definition}{Definition}
\newtheorem{remark}{Remark}
\newcommand{\EqI}{\mathrel{\sim_I}}

\title{P3-T05 Draft/Control Candidate:\\
Invariant-Functor Quotient Relation and Congruence Obligations}
\author{Ontology Formalizer task overlay, RT-20260720-026}
\date{21 July 2026}

\begin{document}
\maketitle

\begin{center}
\fbox{\parbox{0.92\linewidth}{
\textbf{Authority and status.}
This is a \textbf{proposal-only} source-extension construction with status
\texttt{draft/control}.  It is not an adopted source law, not a canonical
ontology edit, not a general discharge of \(\mathrm{EqSrc}\), and not physical
observational equivalence.  Its adoption status is
\texttt{blocked\_adoption\_open\_continuation}.  The construction supplies a
mathematical candidate and bounded controls only; every protected scientific,
Gate Chair, benchmark, and publication decision remains unavailable here.
}}
\end{center}

\section{Question and prospective construction}

The P3-T05 question is whether a general source-side equivalence target can be
formulated without reusing the ad hoc kernel relations of earlier finite
selector families and without defining an invariant from the desired
equivalence classes.  The smallest proposal-neutral answer is the kernel pair,
at object level, of a prospectively fixed invariant functor.

\begin{definition}[Invariant-functor relation]
Let \(\mathcal C\) and \(\mathcal D\) be categories and let
\(I:\mathcal C\to\mathcal D\) be a functor fixed before any pair of source
objects is evaluated.  Work either with a small \(\mathcal C\), or with a
fixed universe-small object collection inside a locally small category.  For
\(A,B\in\operatorname{Ob}(\mathcal C)\), define
\begin{equation}
  A\EqI B
  \quad\Longleftrightarrow\quad
  I(A)\cong I(B)\quad\text{in }\mathcal D .
  \label{eq:kernel-relation}
\end{equation}
The inputs are source objects, source morphisms, the prospectively declared
functor \(I\), and isomorphisms internal to \(\mathcal D\).  Desired class
labels, a target atlas, a target metric, benchmark success, and candidate
reconstruction output are excluded from the definition.
\end{definition}

The word ``invariant'' here has no independent authority.  It abbreviates a
declared functor whose source typing, target neutrality, and admissible domain
must be justified separately.  Equation~\eqref{eq:kernel-relation} would be
circular if \(I\) were built by assigning the same value exactly to objects
already desired to be equivalent.  Prospectivity is therefore part of the
construction, not an optional provenance note.

\section{Exact conditional results}

\begin{theorem}[\texttt{IFQ1\_KERNEL\_EQUIVALENCE}]
For any functor \(I:\mathcal C\to\mathcal D\), the relation \(\EqI\) is an
equivalence relation on \(\operatorname{Ob}(\mathcal C)\).
\end{theorem}
\begin{proof}
For every object \(A\), the identity \(1_{I(A)}\) proves \(A\EqI A\), so the
relation is reflexive.  If \(\phi:I(A)\to I(B)\) is an isomorphism, then
\(\phi^{-1}\) proves \(B\EqI A\), so it is symmetric.  If
\(\phi:I(A)\to I(B)\) and \(\psi:I(B)\to I(C)\) are isomorphisms, then
\(\psi\circ\phi\) proves \(A\EqI C\), so it is transitive.
\end{proof}

\begin{theorem}[\texttt{IFQ2\_MONOIDAL\_COMPOSITION\_CONGRUENCE}]
Suppose \(\mathcal C\) and \(\mathcal D\) have declared binary operations
\(\otimes\) and \(\boxtimes\), the latter is functorial in both arguments, and
there is a natural comparison isomorphism
\[
  \mu_{A,B}:I(A\otimes B)\xrightarrow{\cong} I(A)\boxtimes I(B).
\]
Then \(A\EqI A'\) and \(B\EqI B'\) imply
\(A\otimes B\EqI A'\otimes B'\).
\end{theorem}
\begin{proof}
Choose isomorphisms \(\phi:I(A)\to I(A')\) and
\(\psi:I(B)\to I(B')\).  Bifunctoriality makes
\(\phi\boxtimes\psi\) an isomorphism.  The composite
\[
 I(A\otimes B)
 \xrightarrow{\mu_{A,B}}
 I(A)\boxtimes I(B)
 \xrightarrow{\phi\boxtimes\psi}
 I(A')\boxtimes I(B')
 \xrightarrow{\mu_{A',B'}^{-1}}
 I(A'\otimes B')
\]
is an isomorphism, which is the required congruence witness.
\end{proof}

The theorem proves preservation under the declared composition operation.  It
does not prove cancellation or reflection: those require independent
properties of \(\boxtimes\) and \(I\).

\begin{theorem}[\texttt{IFQ3\_LOCALIZATION\_PRESERVATION}]
Let \(W\) be a declared class of arrows of \(\mathcal C\), let
\(L:\mathcal C\to\mathcal C[W^{-1}]\) be a localization, and suppose
\(I=\bar I\circ L\) up to natural isomorphism.  If \(L(A)\cong L(B)\), then
\(A\EqI B\).  The converse follows only if \(\bar I\) reflects isomorphisms
on the relevant image.
\end{theorem}
\begin{proof}
A functor sends an isomorphism \(L(A)\cong L(B)\) to an isomorphism
\(\bar I L(A)\cong\bar I L(B)\), and the factorization comparisons identify
these objects with \(I(A)\) and \(I(B)\).  Conversely, an isomorphism between
\(\bar I L(A)\) and \(\bar I L(B)\) implies one between \(L(A)\) and
\(L(B)\) precisely when the stated isomorphism-reflection property holds.
\end{proof}

\begin{theorem}[\texttt{IFQ4\_COARSE\_GRAINING\_PRESERVATION}]
Let \(Q:\mathcal C\to\bar{\mathcal C}\) be a prospectively declared
coarse-graining functor and suppose there is a functor
\(\bar I:\bar{\mathcal C}\to\mathcal D\) with a natural isomorphism
\(I\cong\bar I\circ Q\).  Then \(A\EqI B\) is equivalent to
\(\bar I(QA)\cong\bar I(QB)\), so the relation is preserved by this declared
factorizing coarse graining.  No such conclusion follows for an arbitrary
coarse-graining operation.  Reflection of isomorphism in
\(\bar{\mathcal C}\) requires \(\bar I\) to reflect isomorphisms.
\end{theorem}
\begin{proof}
Conjugating an isomorphism \(I(A)\cong I(B)\) by the two components of the
natural comparison gives an isomorphism
\(\bar I(QA)\cong\bar I(QB)\), and the inverse conjugation gives the reverse
implication.  Recovering \(QA\cong QB\) from that image isomorphism is the
separate reflection condition.
\end{proof}

\begin{theorem}[\texttt{IFQ5\_VARIATION\_NATURALITY}]
Let \(V:\mathcal C\to\mathcal C\) be an admissible source variation and let
\(\alpha:I\circ V\Rightarrow I\) be a natural isomorphism.  Then
\[
  A\EqI B\quad\Longleftrightarrow\quad V(A)\EqI V(B).
\]
\end{theorem}
\begin{proof}
Conjugate any \(I(A)\cong I(B)\) by
\(\alpha_A:I(VA)\cong I(A)\) and
\(\alpha_B:I(VB)\cong I(B)\) to obtain an isomorphism
\(I(VA)\cong I(VB)\).  Conjugating with the inverse components proves the
reverse implication.
\end{proof}

More generally, if \(V:\mathcal C\to\mathcal C'\),
\(K:\mathcal D\to\mathcal D'\), and
\(I':\mathcal C'\to\mathcal D'\) admit a natural isomorphism
\(I'V\cong KI\), then \(\EqI\) is preserved by \(V\).  Reflection requires
\(K\) to reflect isomorphisms on the relevant invariant images.  The displayed
theorem is the strict same-invariant special case \(I'=I\) and
\(K=1_{\mathcal D}\).

\begin{proposition}[\texttt{IFQ6\_PHYSICAL\_BRIDGE\_GUARD}]
The preceding category-theoretic results do not establish physical
observational equivalence, physical gauge redundancy, dynamical
interchangeability, or empirical indistinguishability.
\end{proposition}
\begin{proof}
The hypotheses name only categories, functors, comparison isomorphisms, and
source operations.  None supplies an observation map, measurement semantics,
matter response, dynamical law, empirical tolerance, or proof that the
structural arrows are physical redundancies.  Those conclusions therefore do
not follow from the stated premises.
\end{proof}

\section{Finite graph-homology witness model}

The machine-checkable instantiation regards the finite undirected simple graphs
listed in the companion JSON artifact as finite one-dimensional CW complexes
with declared cellular maps.  Integral homology gives the actual functor
\[
  H=(H_0(-;\mathbb Z),H_1(-;\mathbb Z)):
  \mathcal G_{\mathrm{fin}}\longrightarrow
  \mathbf{Ab}_{\mathrm{fg}}\times\mathbf{Ab}_{\mathrm{fg}}.
\]
For these graphs its object-isomorphism class is computed by
\[
 H(G)=(\beta_0(G),\beta_1(G)),\qquad
 \beta_1(G)=|E(G)|-|V(G)|+\beta_0(G),
\]
where the displayed pair records the ranks of the two free abelian groups.
Thus \(G\sim_HG'\) exactly when their Betti pairs agree.  This model is a
bounded witness to the abstract contracts; it is not the proposed physical
substrate invariant.

\begin{center}
\begin{tabular}{>{\raggedright\arraybackslash}p{0.25\linewidth}
                >{\raggedright\arraybackslash}p{0.22\linewidth}
                >{\raggedright\arraybackslash}p{0.41\linewidth}}
\toprule
Control & Exact result & Logical use \\
\midrule
Triangle versus square & \((1,1)=(1,1)\) & Positive pair: nonisomorphic source
objects may have equal invariant images. \\
Triangle versus three-vertex path & \((1,1)\ne(1,0)\) & Negative pair: the
invariant separates a cycle from a tree. \\
One edge versus two disjoint edges & \((1,0)\ne(2,0)\) & Negative pair:
component count remains visible. \\
Disjoint-union context & both become \((2,1)\) & Composition congruence in
the additive finite model. \\
Cycle-edge subdivision & \((1,1)\) is preserved & Declared localization
witness. \\
Pendant-bridge contraction & \((1,1)\) is preserved & Declared exact
coarse-graining witness. \\
Graph relabeling & \((1,1)\) is preserved & Variation naturality witness. \\
Cycle-edge deletion or cycle collapse & \(\beta_1\) changes & Countercontrols:
arbitrary coarse graining is not exact. \\
\bottomrule
\end{tabular}
\end{center}

The finite model also illustrates why preservation and reflection must be
separated.  A coarse map may forget a cycle and merge previously distinct
invariant images; equality after information loss does not reconstruct
equality before it.  The abstract theorems therefore never infer reflection
from the word ``coarse graining.''

\section{Relationship to P3-T03 and P3-T04}

P3-T03 established a conditional factorization result for observables that
are constant on declared structural orbits and supplied countermodels for raw
representative-sensitive interfaces.  P3-T04 established finite probability
and dynamics guards showing that invariant laws or equivariant evolution do
not freely manufacture a non-fixed representative.  The present construction
does not modify either result.  It shifts the target from choosing one
representative to comparing source objects through a prospectively declared
functor.  The earlier results motivate operation and interpretation guards but
do not derive \(\mathcal C\), \(I\), or the admissible operations used here.

\section{Ontology, interpretation, and falsifiers}

The current canonical ontology does not derive the source category
\(\mathcal C\), the invariant category \(\mathcal D\), the functor \(I\), the
localizing class \(W\), the composition comparison \(\mu\), the coarse functor
\(Q\), or the variation class \(V\).  Hence current adoption is blocked while
same-milestone research continuation remains open.  This is
underdetermination, not a theorem of impossibility.

The candidate must be rejected or revised on its declared domain if any of the
following occurs:
\begin{enumerate}
  \item \(I\) cannot be defined without desired EqSrc labels, target-atlas
  data, or benchmark outcomes;
  \item a claimed composition operation lacks a functorial comparison
  isomorphism or yields an explicit congruence counterexample;
  \item a claimed localizing arrow is not inverted by \(I\);
  \item a claimed exact coarse graining does not admit the stated
  factorization, as in the cycle-deletion controls;
  \item an admissible variation lacks the required natural comparison; or
  \item a physical interpretation is asserted without independently derived
  observation, matter-response, or gauge-redundancy structure.
\end{enumerate}

\section{Distance-to-GR and continuation disposition}

This packet adds mathematical payload but changes no Distance-to-GR ledger
entry.  The active milestone remains \texttt{source\_equivalence\_eqsrc}; the
burden remains to replace ad hoc kernel relations with a general
source-equivalence target that is source-derived, audited, stress-tested, and
lawfully adopted.  \(M_{\mathrm{src}}\), \(g_{\mathrm{eff}}\), matter
coupling, the Einstein equations, exact-GR recovery, and benchmark promotion
remain downstream and unestablished.

P3-T05 can complete only as a proposal-only conditional relation packet with
finite witnesses.  The dependency-ready next item is P3-T06, an audit for
circularity, imported ontology, and hidden target dependence.  P3-T06 is
selected but not executed here.

\end{document}
